\documentclass{beamer}
\usetheme{Madrid}

\usepackage[utf8]{inputenc} % solo si usas pdflatex
\usepackage[T1]{fontenc}    % solo si usas pdflatex

\usepackage{amsmath,amssymb}
\usepackage{microtype}
\usepackage{cancel}         % si lo usas
\usepackage{fontawesome}    % si lo usas
% \usepackage{xypic}        % solo si lo usas
% \usepackage{multicol}     % normalmente no; usar columns de beamer
% \usepackage{amsthm}       % solo si defines teoremas propios
% \usepackage{scrextend}    % solo si lo usas

\usepackage{graphicx}
\graphicspath{{images/}}

\hypersetup{pdfborder={0 0 0}}


\title{Introducción a Hadoop}
\author{Rodrigo Blanco Betes. Alfa Formación}
\date{Febrero del 2026}

\begin{document}

\begin{frame}
\titlepage
\end{frame}

\begin{frame}
\tableofcontents
\end{frame}

%-------------------------------------------------

\section{Introducción a Big Data}

\begin{frame}{Introducción a Big Data}
En los últimos años, el crecimiento exponencial de los datos generados por sistemas informáticos, dispositivos móviles, redes sociales, sensores e infraestructuras digitales ha dado lugar al concepto conocido como \textit{Big Data}.

Este término hace referencia a conjuntos de datos cuyo tamaño, complejidad y velocidad de generación superan la capacidad de las herramientas tradicionales de almacenamiento y procesamiento.
\end{frame}

\begin{frame}{Big Data: las cinco Vs}
El Big Data no se define únicamente por el volumen de información, sino por un conjunto de características conocidas como las \textbf{cinco Vs}:
\begin{itemize}
    \item \textbf{Volumen}: grandes cantidades de datos que pueden alcanzar terabytes o petabytes.
    \item \textbf{Velocidad}: rapidez con la que los datos se generan, transmiten y procesan.
    \item \textbf{Variedad}: diversidad de formatos de datos, tanto estructurados como no estructurados.
    \item \textbf{Veracidad}: calidad y fiabilidad de los datos, que pueden contener errores o inconsistencias.
    \item \textbf{Valor}: capacidad de extraer información útil y conocimiento relevante a partir de los datos.
\end{itemize}
\end{frame}

\begin{frame}{Limitaciones y necesidad de plataformas distribuidas}
Los sistemas tradicionales de bases de datos y procesamiento centralizado presentan importantes limitaciones ante estos retos, especialmente en términos de:
\begin{itemize}
    \item \textbf{Escalabilidad}
    \item \textbf{Coste}
    \item \textbf{Tolerancia a fallos}
\end{itemize}

A medida que el volumen de datos crece, resulta inviable aumentar la capacidad de un único sistema sin incurrir en elevados costes económicos y técnicos.

Ante este escenario surge la necesidad de plataformas capaces de distribuir tanto el almacenamiento como el procesamiento entre múltiples nodos. En este contexto, destaca \textbf{Hadoop}, diseñado para gestionar grandes volúmenes de datos de forma eficiente, escalable y tolerante a fallos.
\end{frame}

%----------------------------

\section{Flujo de los datos y ciclo de vida del dato}

\begin{frame}{Flujo de los datos y ciclo de vida del dato}
El tratamiento de datos en entornos Big Data sigue un conjunto de etapas bien definidas que conforman el \textbf{ciclo de vida del dato}.

Este ciclo describe el recorrido que realizan los datos desde su generación hasta su transformación en información útil para la toma de decisiones.
\end{frame}

\begin{frame}{Objetivo del ciclo de vida del dato}
El objetivo principal del ciclo de vida del dato es \textbf{maximizar el valor} de los datos, garantizando:
\begin{itemize}
    \item Captura correcta y fiable desde múltiples fuentes.
    \item Calidad y consistencia.
    \item Almacenamiento adecuado a gran escala.
    \item Procesamiento eficiente.
    \item Presentación clara de resultados.
\end{itemize}

En entornos Big Data, este ciclo suele ser \textbf{continuo y repetitivo}, debido a la generación constante de nuevos datos. Plataformas como \textbf{Hadoop} permiten gestionarlo de manera distribuida, escalable y tolerante a fallos.
\end{frame}

\begin{frame}{Etapas del ciclo de vida del dato (visión general)}
\begin{enumerate}
    \item \textbf{Recolección}
    \item \textbf{Ingestión}
    \item \textbf{Preparación}
    \item \textbf{Procesamiento}
    \item \textbf{Presentación}
\end{enumerate}
\end{frame}

\begin{frame}{Recolección e ingestión de datos}
\textbf{Recolección:} captura inicial desde fuentes heterogéneas:
\begin{itemize}
    \item Sensores IoT, aplicaciones web, redes sociales, sistemas empresariales.
    \item Bases de datos transaccionales, archivos y \textit{logs} de servidores.
    \item Formatos variados: texto, JSON, XML, imágenes, binario.
    \item Generación continua (sensores/logs) o puntual (exportaciones periódicas).
\end{itemize}

\medskip
\textbf{Ingestión:} traslado hacia el sistema de almacenamiento/procesamiento (p.\ ej., Hadoop):
\begin{itemize}
    \item Modo \textbf{batch} (por bloques) o (casi) \textbf{tiempo real} (continuo).
    \item Debe manejar grandes volúmenes y picos de carga, asegurando consistencia y disponibilidad.
\end{itemize}
\end{frame}

\begin{frame}{Preparación y procesamiento de los datos}
\textbf{Preparación:} mejora de la calidad antes del análisis:
\begin{itemize}
    \item Limpieza de errores, eliminación de duplicados.
    \item Gestión de valores nulos.
    \item Transformación y unificación de formatos/codificaciones.
\end{itemize}
Una mala calidad de datos puede afectar negativamente a los resultados, incluso con herramientas avanzadas.

\medskip
\textbf{Procesamiento:} análisis distribuido (p.\ ej., en un clúster Hadoop):
\begin{itemize}
    \item Estadística, agregaciones, búsqueda de patrones.
    \item Clasificación, análisis de tendencias.
    \item Orientado a histórico o a resultados agregados a gran escala.
\end{itemize}
Objetivo: extraer información relevante y generar valor para la organización.
\end{frame}

\begin{frame}{Presentación de los datos}
Fase final centrada en mostrar resultados de forma comprensible:
\begin{itemize}
    \item Informes, gráficos, paneles de control (\textit{dashboards}), visualizaciones interactivas.
    \item Facilita interpretación e identificación de patrones, tendencias o anomalías.
    \item Apoya la toma de decisiones estratégicas (rendimiento, comportamiento de usuarios, KPIs, etc.).
\end{itemize}

Esta fase cierra el ciclo, convirtiendo resultados en información accionable.
\end{frame}

%------------------------------------------

\section{Ejemplo de análisis detallado del ciclo de vida del dato}

\begin{frame}{Ejemplo: análisis detallado del ciclo de vida del dato}
Para ilustrar de forma práctica el ciclo de vida del dato en un entorno Big Data, se presenta un ejemplo que recorre todas las fases:
\begin{itemize}
    \item Identificación del dato en su origen.
    \item Ingesta y organización en zonas de datos.
    \item Preparación y transformaciones.
    \item Cruce y enriquecimiento.
    \item Visualización y KPIs.
\end{itemize}
Este recorrido permite comprender cómo los datos evolucionan a lo largo del proceso y cómo se les añade valor progresivamente.
\end{frame}

\begin{frame}{Fuentes e identificación del dato (e-commerce)}
Los datos proceden de una plataforma digital de comercio electrónico. Fuentes principales:
\begin{itemize}
    \item Registros de acceso de usuarios (\textit{logs}).
    \item Transacciones de compra.
    \item Información básica de productos.
\end{itemize}

\medskip
Naturaleza y formato de los datos:
\begin{itemize}
    \item \textbf{Logs de acceso}: semiestructurados en texto.
    \item \textbf{Transacciones}: estructurados en BBDD relacionales.
    \item \textbf{Productos}: formatos como JSON o CSV.
    \item Generación \textbf{continua} (logs) y \textbf{periódica} (cargas batch).
\end{itemize}
\end{frame}

\begin{frame}{Ingesta y zonas de residencia del dato}
Tras identificar las fuentes, los datos se ingieren en el sistema Big Data:
\begin{itemize}
    \item \textbf{Logs}: incorporación continua.
    \item \textbf{Transacciones y productos}: importación \textbf{batch}.
\end{itemize}

\medskip
Organización por zonas:
\begin{itemize}
    \item \textbf{\textit{Landing zone}}: conservación inicial en formato original.
    \item \textbf{\textit{Raw zone}}: copia íntegra del dato original para trazabilidad y reprocesamientos.
\end{itemize}
\end{frame}

\begin{frame}{Preparación, transformaciones y paso a zona procesada}
En la zona en bruto, se realizan tareas de preparación:
\begin{itemize}
    \item Eliminación de duplicados.
    \item Corrección de errores de formato.
    \item Normalización de fechas, identificadores y valores numéricos.
\end{itemize}

\medskip
Tras las transformaciones:
\begin{itemize}
    \item Los datos se distribuyen en la \textbf{\textit{processed zone}}.
    \item Se almacenan en formatos optimizados para el análisis.
\end{itemize}

La separación por zonas facilita el control del flujo y mejora la trazabilidad.
\end{frame}

\begin{frame}{Cruce de información y enriquecimiento}
Una vez preparados, los datos se cruzan y enriquecen:
\begin{itemize}
    \item Combinación de \textit{logs} de acceso con transacciones para analizar patrones de comportamiento.
    \item Uso de la información de productos para contextualizar compras y segmentar resultados.
\end{itemize}

\medskip
Resultado:
\begin{itemize}
    \item Conjuntos de datos analíticos integrados desde múltiples fuentes.
    \item Mayor valor informativo y visión más completa del negocio.
\end{itemize}
\end{frame}

\begin{frame}{Visualización, cuadros de mando y KPIs}
En la fase final, los resultados se presentan mediante herramientas de visualización:
\begin{itemize}
    \item Construcción de cuadros de mando (\textit{dashboards}).
    \item Indicadores clave (KPIs), por ejemplo:
    \begin{itemize}
        \item Número de usuarios activos.
        \item Volumen de ventas.
        \item Tasa de conversión.
    \end{itemize}
\end{itemize}

Estas visualizaciones permiten interpretar resultados y tomar decisiones basadas en datos, completando el ciclo: de información bruta a conocimiento útil.
\end{frame}

%-----------------------------------------

\section{Hadoop}

\begin{frame}{¿Qué es Hadoop?}
Apache Hadoop es un framework de software de código abierto diseñado para el almacenamiento y procesamiento distribuido de grandes volúmenes de datos. Permite utilizar clústeres formados por hardware convencional para tratar datos de forma eficiente, escalable y tolerante a fallos.

Se basa en la \textbf{escalabilidad horizontal}, aumentando la capacidad del sistema mediante la incorporación de nuevos nodos al clúster, en lugar de potenciar un único servidor. Este enfoque reduce costes y mejora la flexibilidad frente a grandes cargas de trabajo.
\end{frame}

\begin{frame}{Objetivos y relevancia de Hadoop}
Entre los objetivos fundamentales de Hadoop destacan:
\begin{itemize}
    \item Procesar grandes conjuntos de datos de forma distribuida.
    \item Garantizar tolerancia a fallos mediante replicación de datos.
    \item Ofrecer alta disponibilidad y fiabilidad.
    \item Acercar el procesamiento a los datos para minimizar el tráfico de red.
\end{itemize}

Gracias a estas características, Hadoop se ha consolidado como una de las tecnologías base del ecosistema Big Data, con amplio uso en entornos académicos y empresariales.
\end{frame}

\begin{frame}{Arquitectura general de Hadoop}
La arquitectura de Hadoop sigue un modelo distribuido basado en un clúster de nodos con estructura \textit{master--slave}. Cada nodo desempeña un rol específico que permite separar:
\begin{itemize}
    \item Gestión del sistema.
    \item Almacenamiento de datos.
    \item Procesamiento distribuido.
\end{itemize}

Se distinguen dos tipos principales de nodos:
\begin{itemize}
    \item \textbf{Nodos maestros}: gestión global y asignación de recursos.
    \item \textbf{Nodos trabajadores}: almacenamiento de datos y ejecución de tareas.
\end{itemize}
\end{frame}

\begin{frame}{Módulos principales de Hadoop}
Hadoop está compuesto por varios módulos que cooperan entre sí:
\begin{itemize}
    \item \textbf{HDFS}: sistema de archivos distribuido para el almacenamiento.
    \item \textbf{YARN}: gestión de recursos y planificación de tareas.
    \item \textbf{MapReduce}: modelo de programación para el procesamiento distribuido.
\end{itemize}

La comunicación entre estos componentes permite detectar y gestionar fallos automáticamente, redistribuyendo tareas y garantizando la continuidad del sistema.
\end{frame}

\begin{frame}{División de datos en bloques (HDFS)}
HDFS divide los archivos de gran tamaño en bloques de tamaño fijo que se distribuyen entre los nodos del clúster:
\begin{itemize}
    \item Tamaño de bloque habitual: \textbf{128 MB o 256 MB}.
    \item Reducción de operaciones de E/S y mejora del rendimiento secuencial.
    \item Replicación de bloques en distintos nodos para tolerancia a fallos.
\end{itemize}

El \textbf{NameNode} gestiona la ubicación y replicación de los bloques, mientras que la división en bloques facilita el procesamiento paralelo y la escalabilidad del sistema.
\end{frame}

%----------------------------

\section{Componentes de Hadoop}

\begin{frame}{Componentes principales de Hadoop}
Hadoop se basa en una arquitectura modular cuyos componentes trabajan de forma coordinada para ofrecer:
\begin{itemize}
    \item Almacenamiento distribuido.
    \item Gestión eficiente de recursos.
    \item Procesamiento paralelo de grandes volúmenes de datos.
\end{itemize}

Los componentes clave son:
\begin{itemize}
    \item HDFS
    \item YARN
    \item MapReduce
\end{itemize}
\end{frame}

\begin{frame}{HDFS y almacenamiento distribuido}
HDFS (Hadoop Distributed File System) es el sistema de archivos distribuido de Hadoop. Su objetivo es almacenar grandes volúmenes de datos de forma fiable y escalable.

Desde el punto de vista arquitectónico, separa:
\begin{itemize}
    \item Gestión de metadatos.
    \item Almacenamiento físico de los datos.
\end{itemize}
\end{frame}

\begin{frame}{NameNode y DataNodes}
El \textbf{NameNode} gestiona los metadatos del sistema de archivos:
\begin{itemize}
    \item Estructura de directorios y archivos.
    \item Ubicación de los bloques.
    \item Permisos y políticas de acceso.
    \item Replicación para tolerancia a fallos.
\end{itemize}

Los \textbf{DataNodes} (WorkerNodes) almacenan físicamente los bloques y:
\begin{itemize}
    \item Atienden lecturas y escrituras.
    \item Informan periódicamente de su estado al NameNode.
\end{itemize}
\end{frame}

\begin{frame}{YARN y la gestión de recursos}
YARN es el componente encargado de la gestión de recursos y la planificación de aplicaciones en el clúster.

Su objetivo principal es asignar de forma eficiente recursos como CPU y memoria, permitiendo la ejecución simultánea de múltiples aplicaciones y frameworks de procesamiento.
\end{frame}

\begin{frame}{Componentes internos de YARN}
YARN se compone de varios elementos clave:
\begin{itemize}
    \item \textbf{ResourceManager}: gestión global de recursos y planificación.
    \item \textbf{NodeManager}: control de recursos y contenedores en cada nodo.
    \item \textbf{ApplicationMaster}: gestión del ciclo de vida de cada aplicación.
    \item \textbf{Containers}: unidad básica de asignación de CPU y memoria.
\end{itemize}
\end{frame}

\begin{frame}{MapReduce como modelo de procesamiento}
MapReduce es un modelo de programación para procesar grandes volúmenes de datos de forma distribuida y paralela.

Se basa en:
\begin{itemize}
    \item \textbf{Map}: generación de pares clave--valor intermedios.
    \item \textbf{Shuffle y Sort}: agrupación de pares por clave.
    \item \textbf{Reduce}: combinación de valores para obtener resultados finales.
\end{itemize}
\end{frame}

\begin{frame}{Usos y características de MapReduce}
MapReduce se utiliza habitualmente para:
\begin{itemize}
    \item Cálculo de estadísticas agregadas.
    \item Análisis de logs y patrones de uso.
    \item Procesamiento de datos históricos.
\end{itemize}

Aunque es robusto y tolerante a fallos, su enfoque batch y su latencia han motivado el uso de otros modelos más flexibles en ciertos escenarios, si bien sigue siendo un pilar fundamental de la arquitectura Hadoop.
\end{frame}

%---------------------------------

\section{Ecosistema Hadoop}

\begin{frame}{Ecosistema Hadoop}
Hadoop no es una herramienta aislada, sino el núcleo de un amplio \textbf{ecosistema de proyectos} diseñados para cubrir las distintas necesidades del tratamiento de datos en entornos Big Data.

Este ecosistema amplía las capacidades básicas de:
\begin{itemize}
    \item Almacenamiento distribuido.
    \item Procesamiento paralelo.
    \item Gestión eficiente de recursos.
\end{itemize}

Proporciona herramientas especializadas para consulta, análisis, ingesta, bases de datos distribuidas y orquestación de procesos.
\end{frame}

\begin{frame}{Integración sobre HDFS y YARN}
Las herramientas del ecosistema Hadoop se integran principalmente sobre:
\begin{itemize}
    \item \textbf{HDFS}: almacenamiento distribuido y tolerante a fallos.
    \item \textbf{YARN}: gestión centralizada de recursos y ejecución de aplicaciones.
\end{itemize}

Esta integración permite ofrecer soluciones:
\begin{itemize}
    \item Escalables.
    \item Distribuidas.
    \item Altamente disponibles.
\end{itemize}
\end{frame}

\begin{frame}{Hive: consultas tipo SQL}
\textbf{Apache Hive} es una herramienta de análisis que permite consultar datos almacenados en HDFS mediante un lenguaje similar a SQL (\textit{HiveQL}).

Características principales:
\begin{itemize}
    \item Orientado a usuarios con experiencia en bases de datos relacionales.
    \item Traduce consultas en trabajos de procesamiento distribuido.
    \item Especialmente adecuado para análisis \textbf{batch} y datos históricos.
\end{itemize}

Es habitual su uso en generación de informes y análisis exploratorio.
\end{frame}

\begin{frame}{Pig: flujos de transformación de datos}
\textbf{Apache Pig} proporciona un lenguaje de alto nivel llamado \textit{Pig Latin}, orientado a describir flujos de procesamiento de datos.

Se utiliza principalmente para:
\begin{itemize}
    \item Limpieza y filtrado de datos.
    \item Transformaciones y agregaciones.
    \item Combinación de conjuntos de datos.
\end{itemize}

Actúa como una capa intermedia entre los datos almacenados en HDFS y los procesos analíticos posteriores.
\end{frame}

\begin{frame}{HBase: base de datos NoSQL distribuida}
\textbf{Apache HBase} es una base de datos NoSQL distribuida que se ejecuta sobre HDFS.

A diferencia del enfoque batch de Hadoop, HBase está orientada a:
\begin{itemize}
    \item Acceso rápido y aleatorio a los datos.
    \item Operaciones de lectura y escritura en tiempo casi real.
\end{itemize}

Es especialmente útil en aplicaciones que requieren baja latencia, como monitorización, perfiles de usuario o acceso inmediato a datos semiestructurados.
\end{frame}

\begin{frame}{Sqoop: integración con bases de datos relacionales}
\textbf{Apache Sqoop} permite la transferencia eficiente de datos entre Hadoop y bases de datos relacionales.

Funcionalidades clave:
\begin{itemize}
    \item Importación de datos desde sistemas tradicionales hacia HDFS, Hive o HBase.
    \item Exportación de resultados procesados desde Hadoop a bases de datos externas.
\end{itemize}

Facilita la integración de Hadoop con infraestructuras ya existentes en las organizaciones.
\end{frame}

\begin{frame}{Flume: ingesta de datos en tiempo casi real}
\textbf{Apache Flume} es un sistema distribuido para la recolección y transporte de grandes volúmenes de datos en tiempo casi real.

Uso principal:
\begin{itemize}
    \item Ingesta de logs, eventos y flujos continuos.
    \item Datos generados por aplicaciones y servidores.
\end{itemize}

Está diseñado para ser escalable y tolerante a fallos, asegurando una transferencia fiable hacia HDFS u otros sistemas del ecosistema.
\end{frame}

\begin{frame}{Oozie: orquestación de flujos de trabajo}
\textbf{Apache Oozie} es una herramienta para la gestión y coordinación de flujos de trabajo en Hadoop.

Permite:
\begin{itemize}
    \item Definir secuencias de trabajos de ingesta, procesamiento y análisis.
    \item Programar tareas y gestionar dependencias.
    \item Automatizar procesos complejos de principio a fin.
\end{itemize}
\end{frame}

\begin{frame}{Relación entre las herramientas del ecosistema}
Las herramientas del ecosistema Hadoop se complementan para cubrir todo el ciclo de vida del dato:
\begin{itemize}
    \item \textbf{Flume / Sqoop}: ingesta de datos.
    \item \textbf{HDFS}: almacenamiento distribuido.
    \item \textbf{Pig / Hive}: preparación y análisis.
    \item \textbf{HBase}: acceso rápido a resultados.
    \item \textbf{Oozie}: coordinación y orquestación.
\end{itemize}

Esta integración convierte a Hadoop en una \textbf{plataforma completa} para el tratamiento de grandes volúmenes de datos en entornos distribuidos.
\end{frame}

%----------------------------------------

\section{Consideraciones finales sobre Hadoop}

\begin{frame}{Consideraciones finales sobre Hadoop}
En esta sección se resumen algunos aspectos clave relacionados con el uso de Hadoop:
\begin{itemize}
    \item Principales ventajas y desventajas.
    \item Casos de uso habituales.
    \item Papel actual de Hadoop en el ecosistema Big Data.
\end{itemize}
\end{frame}

\begin{frame}{Ventajas de Hadoop}
Hadoop presenta numerosas ventajas que han favorecido su adopción en entornos Big Data:
\begin{itemize}
    \item \textbf{Escalabilidad horizontal} mediante la incorporación de nuevos nodos.
    \item \textbf{Tolerancia a fallos} gracias a la replicación de datos y ejecución distribuida.
    \item \textbf{Uso de hardware convencional}, reduciendo costes de infraestructura.
    \item Capacidad para procesar \textbf{grandes volúmenes de datos} de forma eficiente.
\end{itemize}
\end{frame}

\begin{frame}{Desventajas de Hadoop}
A pesar de sus ventajas, Hadoop también presenta ciertas limitaciones:
\begin{itemize}
    \item \textbf{Elevada latencia} frente a sistemas orientados al tiempo real.
    \item \textbf{Complejidad} en la instalación, configuración y administración del clúster.
    \item Menor eficiencia con \textbf{conjuntos de datos pequeños}.
    \item \textbf{Curva de aprendizaje pronunciada} para nuevos usuarios.
\end{itemize}
\end{frame}

\begin{frame}{Casos de uso habituales}
Hadoop se utiliza en escenarios donde se requiere el procesamiento y análisis de grandes volúmenes de datos:
\begin{itemize}
    \item Procesamiento y análisis de archivos de registros (\textit{logs}).
    \item Análisis de datos en comercio electrónico, telecomunicaciones y redes sociales.
    \item Sistemas de recomendación y análisis del comportamiento de usuarios.
    \item Procesamiento de datos de sensores y entornos IoT.
\end{itemize}
\end{frame}

\begin{frame}{Hadoop en la actualidad}
Aunque Hadoop fue durante años la tecnología de referencia en Big Data, hoy comparte protagonismo con soluciones más modernas, como Apache Spark, que ofrecen:
\begin{itemize}
    \item Menor latencia.
    \item Mayor flexibilidad en los modelos de procesamiento.
\end{itemize}

No obstante, Hadoop sigue siendo fundamental en muchos entornos empresariales, especialmente como base para:
\begin{itemize}
    \item Almacenamiento distribuido.
    \item Gestión de recursos a gran escala.
\end{itemize}
\end{frame}

\begin{frame}{Conclusión}
Hadoop ha desempeñado un papel clave en la evolución del Big Data, sentando las bases del procesamiento distribuido a gran escala.

Su arquitectura escalable y tolerante a fallos lo convierte en una solución robusta para el tratamiento de grandes volúmenes de datos, manteniendo su relevancia como componente esencial dentro de los ecosistemas modernos de análisis de datos.
\end{frame}


\end{document}
